\documentclass{standalone}
\usepackage{ctex}
\usepackage{tikz}

\usetikzlibrary{
    arrows.meta,
    positioning,
    calc,
    fit,
    backgrounds,
    shapes.geometric
}

\tikzset{
    actor/.style = {
        draw,
        rectangle,
        rounded corners=6pt,
        align=center,
        minimum width=2.7cm,
        minimum height=1.0cm,
        font=\small\bfseries
    },
    panel/.style = {
        draw,
        rectangle,
        rounded corners=6pt,
        align=left,
        inner sep=8pt,
        font=\small,
        text width=7.1cm
    },
    service/.style = {
        draw,
        rectangle,
        rounded corners=6pt,
        align=left,
        inner sep=8pt,
        font=\small,
        text width=7.1cm
    },
    db/.style = {
        draw,
        cylinder,
        shape border rotate=90,
        aspect=0.25,
        align=left,
        minimum width=3.9cm,
        minimum height=6.4cm,
        font=\small
    },
    resource/.style = {
        draw,
        rectangle,
        rounded corners=5pt,
        align=center,
        minimum width=3.5cm,
        minimum height=0.95cm,
        font=\small
    },
    arrow/.style = {
        ->,
        thick,
        >={Stealth[length=3mm,width=2mm]}
    }
}

\begin{document}

\begin{tikzpicture}[
    node distance=1.4cm and 2.4cm
]

%================================================
% 1. 前端应用层：先上下排好两个前端模块
%    admin 位于 frontend 下方，边到边距离为 1.2cm
%================================================

\node[panel] (frontend) {
\textbf{微信小程序前台（frontend）}\\[5pt]
1. 首页：展示站点信息、置顶文章、最新文章\\
2. 归档页：按时间浏览所有公开文章\\
3. 友链页：展示项目技术栈相关官方链接\\
4. 文章详情页：展示标题、封面、标签、正文等
};

\node[panel, below=1.2cm of frontend] (admin) {
\textbf{后台管理端（admin）}\\[5pt]
1. 登录认证\\
2. 文章管理\\
3. 友链管理\\
4. 系统配置\\
5. 邮件日志管理
};

%================================================
% 2. 不可见辅助框 frontgroup
%    这个框包住 frontend 和 admin
%    后面的 backend 就相对这个整体定位
%================================================

\node[
    fit=(frontend)(admin),
    inner sep=0pt,
    outer sep=0pt
] (frontgroup) {};

%================================================
% 3. 用户层
%    用户距离对应前端模块左边界 2.4cm
%================================================

\node[actor, left=2.4cm of frontend] (visitor) {普通用户};
\node[actor, left=2.4cm of admin] (adminuser) {管理员};

%================================================
% 4. 业务服务层
%    backend 距离整个 frontgroup 的右边界 3.2cm
%================================================

\node[service, right=3.2cm of frontgroup] (backend) {
\textbf{后端服务（backend）}\\[5pt]
1. Auth 路由：登录、退出、找回密码\\
2. Post 路由：文章增删改查、公开文章接口\\
3. Friend 路由：友链申请、审核、展示\\
4. Config 路由：站点配置读取与保存\\
5. Upload 路由：图片上传与资源管理\\
6. Email 路由：SMTP 测试、邮件日志查询
};

%================================================
% 5. 数据与资源层
%    database 距离 backend 右边界 3.2cm
%================================================

\node[db, right=3.2cm of backend] (database) {
\textbf{SQLite 数据库}\\[8pt]
\texttt{user}\\[2pt]
\texttt{post}\\[2pt]
\texttt{friend\_link}\\[2pt]
\texttt{config}\\[2pt]
\texttt{email\_log}\\[2pt]
\texttt{email\_verification}
};

\node[resource, above=1.4cm of database] (upload) {
上传资源目录
};

\node[resource, below=1.4cm of database] (mail) {
邮件服务 / 邮件日志
};

%================================================
% 6. 连线
%================================================

\draw[arrow] (visitor.east) -- (frontend.west);
\draw[arrow] (adminuser.east) -- (admin.west);

\draw[arrow] (frontend.east) -- (backend.west);
\draw[arrow] (admin.east) -- (backend.west);

\draw[arrow] (backend.east) -- (database.west);

\draw[arrow] (backend.east) -- ++(1.2cm,0) |- (upload.west);
\draw[arrow] (backend.east) -- ++(1.2cm,0) |- (mail.west);

%================================================
% 7. 分层虚线框
%================================================

\begin{pgfonlayer}{background}

\node[
    draw,
    dashed,
    rounded corners=8pt,
    inner sep=12pt,
    fit=(frontend)(admin),
    label={[font=\bfseries]above:前端应用层}
] {};

\node[
    draw,
    dashed,
    rounded corners=8pt,
    inner sep=12pt,
    fit=(backend),
    label={[font=\bfseries]above:业务服务层}
] {};

\node[
    draw,
    dashed,
    rounded corners=8pt,
    inner sep=12pt,
    fit=(database)(upload)(mail),
    label={[font=\bfseries]above:数据与资源层}
] {};

\end{pgfonlayer}

\end{tikzpicture}

\end{document}
